


\begin{abstract}
Generative user interfaces have become an active topic in human-computer interaction and AI-supported software design because they promise interfaces that adapt to user goals instead of remaining fully static. Yet this promise creates a practical tension for productivity software, where users depend on stable layouts, predictable controls, and durable mental models. This thesis investigates whether ephemeral, task-scoped generative UI components can offer a bounded alternative to full interface regeneration by appearing only when locally useful and disappearing without destabilising the host application.

The thesis develops and tests a framework for ephemeral support through a web-based research artifact. The main individual contributions are the framework operationalisation itself, a constrained component-specification model for temporary support, a validation pipeline for generated interface specifications, a bounded ephemeral overlay approach that preserves interface continuity, and an instrumented comparative study artifact for empirical evaluation. The artifact was tested in a within-subjects user study comparing a baseline interface with an ephemeral condition on a bounded slide-reordering task.

The results show a mixed outcome. The ephemeral condition produced a slightly higher task completion rate and somewhat higher helpfulness ratings, while perceived control remained stable, but it did not improve completion time or overall preference and significantly increased interaction count and perceived intrusiveness. These findings suggest that ephemeral generative support is most credible not as a general replacement for existing interfaces, but as a bounded design strategy for low-risk productivity tasks where lightweight contextual guidance may be useful. The thesis therefore contributes both a concrete design framework and an empirical basis for assessing when temporary generative assistance can add value without disrupting the user experience.
\end{abstract}

Keywords: generative UI, ephemeral interfaces, human-computer interaction, productivity software, design science research
